\documentclass{article}
\usepackage{amsmath,amssymb,amsthm,algorithm,algorithmic}
\usepackage{bm}
\usepackage{hyperref}
\usepackage{enumitem}
\newtheorem{theorem}{Theorem}
\newtheorem{lemma}{Lemma}
\newtheorem{assumption}{Assumption}
\newtheorem{corollary}{Corollary}
\newcommand{\R}{\mathbb{R}}
\newcommand{\C}{\mathcal{C}}
\newcommand{\V}{\mathcal{V}}
\newcommand{\gap}{\mathcal{G}}
\newcommand{\diam}{\mathrm{diam}}
\newcommand{\width}{\delta}
\newcommand{\curv}{C_f}
\begin{document}

%%%%%%%%%%%%%%%%%%%%%%%%%%%%%%%%%%%%%%%%%%%%%%%%%%%%%%%%%%%%
\section*{Away‑Step Frank–Wolfe (AFW)}
%%%%%%%%%%%%%%%%%%%%%%%%%%%%%%%%%%%%%%%%%%%%%%%%%%%%%%%%%%%%

\section*{Mathematical Formulation}
Let $f:\R^d\rightarrow\R$ be convex and $L$–smooth.  
Let $\C\subset\R^d$ be a compact polytope,
\[
\C=\operatorname{conv}(\V),\qquad \V=\{v_1,\dots ,v_m\}
\]
(the set of vertices).  
For any iterate $x_k\in\C$ we maintain an \emph{active set}
\[
S_k\subseteq \V,\qquad
x_k=\sum_{v\in S_k}\alpha_v^{(k)}v,\qquad
\alpha_v^{(k)}\ge 0,\ \sum_{v\in S_k}\alpha_v^{(k)}=1 .
\]

\paragraph{Step 1 – Linear minimization (FW vertex).}
\[
s_k \;:=\; \arg\min_{v\in\V}\,\langle\nabla f(x_k),v\rangle .
\tag{1}
\]

\paragraph{Step 2 – Away vertex.}
\[
v_k \;:=\; \arg\max_{v\in S_k}\,\langle\nabla f(x_k),v\rangle .
\tag{2}
\]

\paragraph{Step 3 – Choose direction.}
\[
d_k^{\text{FW}} := s_k-x_k,\qquad
d_k^{\text{A}}  := x_k-v_k .
\]
If 
\[
\langle\nabla f(x_k),d_k^{\text{FW}}\rangle\le
\langle\nabla f(x_k),d_k^{\text{A}}\rangle,
\]
set $d_k:=d_k^{\text{FW}}$ (a \emph{Frank–Wolfe step}); otherwise $d_k:=d_k^{\text{A}}$ (an \emph{away step}).

\paragraph{Step 4 – Maximum feasible step‑size.}
\[
\gamma_{\max}:=
\begin{cases}
1, & d_k=d_k^{\text{FW}},\\[4pt]
\displaystyle\frac{\alpha_{v_k}^{(k)}}{1-\alpha_{v_k}^{(k)}},
& d_k=d_k^{\text{A}} .
\end{cases}
\tag{3}
\]

\paragraph{Step 5 – Line‑search (exact).}
\[
\gamma_k \;:=\;
\arg\min_{\gamma\in[0,\gamma_{\max}]}
f\!\bigl(x_k+\gamma d_k\bigr).
\tag{4}
\]

\paragraph{Step 6 – Update.}
\[
x_{k+1}=x_k+\gamma_k d_k .
\tag{5}
\]
The coefficients $\{\alpha_v^{(k)}\}$ are updated in the obvious way:
\begin{itemize}
\item FW step: increase $\alpha_{s_k}^{(k)}$ by $\gamma_k$ and decrease all others proportionally;
\item Away step: decrease $\alpha_{v_k}^{(k)}$ by $\gamma_k$; if $\alpha_{v_k}^{(k)}$ becomes zero, remove $v_k$ from $S_k$.
\end{itemize}
The algorithm stops when the \emph{Frank–Wolfe duality gap}
\[
\gap_k:=\langle\nabla f(x_k),x_k-s_k\rangle
\]
falls below a prescribed tolerance.

\section*{Pseudocode}
\begin{algorithm}[H]
\caption{Away‑Step Frank–Wolfe (AFW)}
\begin{algorithmic}[1]
\STATE \textbf{Input:} convex $L$‑smooth $f$, polytope $\C=\operatorname{conv}(\V)$,
tolerance $\varepsilon>0$, maximum iterations $T$.
\STATE Choose $x_0\in\V$, set $S_0=\{x_0\}$, $\alpha_{x_0}^{(0)}=1$.
\FOR{$k=0,1,\dots,T-1$}
    \STATE Compute gradient $g_k:=\nabla f(x_k)$.
    \STATE \textbf{FW vertex:} $s_k\gets\arg\min_{v\in\V}\langle g_k,v\rangle$.
    \STATE \textbf{Away vertex:} $v_k\gets\arg\max_{v\in S_k}\langle g_k,v\rangle$.
    \STATE $d_k^{\text{FW}}:=s_k-x_k$, $d_k^{\text{A}}:=x_k-v_k$.
    \IF{$\langle g_k,d_k^{\text{FW}}\rangle\le\langle g_k,d_k^{\text{A}}\rangle$}
        \STATE $d_k\gets d_k^{\text{FW}}$, $\gamma_{\max}\gets 1$.
    \ELSE
        \STATE $d_k\gets d_k^{\text{A}}$, 
        $\displaystyle\gamma_{\max}\gets\frac{\alpha_{v_k}^{(k)}}{1-\alpha_{v_k}^{(k)}}$.
    \ENDIF
    \STATE \textbf{Exact line‑search:}
    \[
    \gamma_k\gets\arg\min_{\gamma\in[0,\gamma_{\max}]}
    f(x_k+\gamma d_k).
    \]
    \STATE $x_{k+1}\gets x_k+\gamma_k d_k$.
    \STATE Update coefficients $\{\alpha_v^{(k)}\}$ and active set $S_{k+1}$.
    \STATE Compute duality gap $\gap_k:=\langle g_k,x_k-s_k\rangle$.
    \IF{$\gap_k\le\varepsilon$}
        \RETURN $x_{k+1}$.
    \ENDIF
\ENDFOR
\RETURN $x_T$.
\end{algorithmic}
\end{algorithm}

\section*{Dry Run Example}
Consider the one‑dimensional problem
\[
f(w)=(w-3)^2,\qquad \C=[0,5]=\operatorname{conv}\{0,5\},
\]
with $w_0=0$ and exact line‑search (no fixed $\eta$).  
We perform three AFW iterations.

\begin{center}
\begin{tabular}{c|c|c|c|c}
\textbf{Iter.} & $w_k$ & $\nabla f(w_k)=2(w_k-3)$ & $s_k$ (FW vertex) & $v_k$ (Away vertex) \\ \hline
0 & $0$ & $-6$ & $5$ (since $-6\cdot5<-6\cdot0$) & $0$ (only active vertex)\\
1 & $3$ & $0$  & $0$ (any vertex gives $0$, pick $0$) & $5$ (both active, pick $5$)\\
2 & $3$ & $0$  & $0$ & $5$ (same as previous)\\
\end{tabular}
\end{center}

\paragraph{Iteration 0 (FW step).}
\[
d_0 = s_0-w_0 = 5-0 = 5,\qquad
\gamma_{\max}=1.
\]
Exact line‑search:
\[
\min_{\gamma\in[0,1]} f(0+5\gamma) = (5\gamma-3)^2
\Longrightarrow \gamma_0=\frac{3}{5}=0.6 .
\]
\[
w_1 = 0 + 0.6\cdot5 = 3,\qquad
f(w_1)=0.
\]

\paragraph{Iteration 1 (FW direction chosen).}
\[
d_1^{\text{FW}} = s_1-w_1 = 0-3 = -3,\qquad
d_1^{\text{A}} = w_1-v_1 = 3-5 = -2.
\]
Both directional derivatives are $0$ because $\nabla f(w_1)=0$, so we keep the FW direction.
Exact line‑search on $f(3+\gamma(-3))=( -3\gamma)^2$ yields $\gamma_1=0$.
Thus $w_2=w_1=3$, $f(w_2)=0$.

\paragraph{Iteration 2.}
Identical to iteration 1; the algorithm stays at the optimum $w^\star=3$.

The table below records the objective values:

\[
\begin{array}{c|c}
k & f(w_k)\\ \hline
0 & 9\\
1 & 0\\
2 & 0\\
\end{array}
\]

The AFW algorithm has reached the optimum after a single away‑free Frank–Wolfe step.

%%%%%%%%%%%%%%%%%%%%%%%%%%%%%%%%%%%%%%%%%%%%%%%%%%%%%%%%%%%%
\section*{Assumptions}
%%%%%%%%%%%%%%%%%%%%%%%%%%%%%%%%%%%%%%%%%%%%%%%%%%%%%%%%%%%%
\begin{assumption}[Convexity and smoothness]\label{ass:convex_smooth}
$f:\R^d\rightarrow\R$ is convex and $L$‑smooth on $\C$, i.e.
\[
\|\nabla f(x)-\nabla f(y)\|\le L\|x-y\|,\qquad\forall x,y\in\C .
\tag{A1}
\]
\end{assumption}
\begin{assumption}[Polytope geometry]\label{ass:polytope}
$\C=\operatorname{conv}(\V)$ is a compact polytope with
\[
\diam(\C):=\max_{u,v\in\C}\|u-v\|\;<\infty .
\tag{A2}
\]
Define the \emph{pyramidal width} $\delta>0$ of $\C$ (see \cite{lacoste2015linear}) as
\[
\delta:=\min_{\substack{x\in\C,\;g\neq0\\
s\in\V,\;v\in S(x)}}
\frac{\langle g,\,x-s\rangle}{\|g\|\,\|x-s\|}
\quad\text{s.t. } \langle g,\,x-v\rangle\ge\langle g,\,x-s\rangle ,
\tag{A3}
\]
where $S(x)$ denotes the set of vertices with positive weight in a representation of $x$.
\end{assumption}
\begin{assumption}[Bounded curvature]\label{ass:curvature}
Define the curvature constant
\[
\curv:=\sup_{\substack{x\in\C,\;s\in\V\\
\gamma\in[0,1]}}
\frac{2}{\gamma^{2}}\Bigl(
f\bigl(x+\gamma(s-x)\bigr)-f(x)-\gamma\langle\nabla f(x),s-x\rangle
\Bigr).
\tag{A4}
\]
Because of $L$‑smoothness, $\curv\le L\diam(\C)^{2}$.
\end{assumption}

%%%%%%%%%%%%%%%%%%%%%%%%%%%%%%%%%%%%%%%%%%%%%%%%%%%%%%%%%%%%
\section*{Main Convergence Theorem}
%%%%%%%%%%%%%%%%%%%%%%%%%%%%%%%%%%%%%%%%%%%%%%%%%%%%%%%%%%%%
\begin{theorem}[Linear convergence of AFW]\label{thm:linear}
Let Assumptions \ref{ass:convex_smooth}–\ref{ass:curvature} hold.
Denote $f^\star:=\min_{x\in\C}f(x)$ and the primal error $h_k:=f(x_k)-f^\star$.
Define the constant
\[
\rho:=\frac{\delta^{2}}{L\,\diam(\C)^{2}}\;\in(0,1].
\tag{7}
\]
Then for every iteration $k\ge0$,
\[
h_{k+1}\;\le\;(1-\rho)\,h_{k}.
\tag{8}
\]
Consequently,
\[
h_{k}\;\le\;(1-\rho)^{k}\,h_{0},
\qquad
\gap_k\;\le\;2L\,\diam(\C)^{2}\,(1-\rho)^{k}\,h_{0}.
\tag{9}
\]
\end{theorem}
\begin{proof}[Sketch of the proof]
The proof follows three lemmas:
\begin{enumerate}[label={\bf Lemma \arabic*}:, leftmargin=*]
\item A descent lemma linking the curvature constant $\curv$ to the actual decrease (Lemma~\ref{lem:descent}).
\item A geometric lemma that lower‑bounds the duality gap $\gap_k$ by $\delta\sqrt{2Lh_k}$ (Lemma~\ref{lem:gap}).
\item Combining the two lemmas yields the linear contraction (8).
\end{enumerate}
Full details are given in Section~\ref{sec:proof}.
\end{proof}

\section*{Theoretical Properties}
\begin{itemize}
\item \textbf{Stability.} The step‑size $\gamma_k$ is always bounded by $\gamma_{\max}\le1$, guaranteeing that iterates stay inside $\C$.
\item \textbf{Hyper‑parameter sensitivity.} AFW does not require a learning‑rate schedule; the only algorithmic parameters are the line‑search rule (exact vs.~Armijo) and the tolerance $\varepsilon$. The linear rate (7) depends solely on geometric constants $\delta$, $L$, and $\diam(\C)$.
\item \textbf{Comparison.} Classical Frank–Wolfe attains a sublinear $O(1/k)$ rate under the same assumptions. Adding away steps (or equivalently the ``pairwise'' variant) upgrades the rate to linear whenever $\delta>0$, i.e. for all polytopes.
\end{itemize}

\section*{Discussion}
The constant $\rho$ in \eqref{7} encapsulates the interaction between smoothness ($L$) and the polytope’s geometry ($\delta$, $\diam$).  
For a simplex in $\R^{d}$, $\delta=2/\sqrt{d(d+1)}$ and $\diam=\sqrt{2}$, yielding an explicit $\rho\sim 1/(Ld)$.  
If $\C$ is a line segment, $\delta=1$ and the bound reduces to $\rho=1/(L\diam^{2})$, which matches the exact contraction observed in the dry‑run example.

%%%%%%%%%%%%%%%%%%%%%%%%%%%%%%%%%%%%%%%%%%%%%%%%%%%%%%%%%%%%
\section*{Preliminaries (Definitions and Lemmas)}\label{sec:prelim}
%%%%%%%%%%%%%%%%%%%%%%%%%%%%%%%%%%%%%%%%%%%%%%%%%%%%%%%%%%%%
\begin{lemma}[Descent Lemma]\label{lem:descent}
For any $x\in\C$, $s\in\V$, and $\gamma\in[0,1]$,
\[
f\bigl(x+\gamma(s-x)\bigr)\le
f(x)+\gamma\langle\nabla f(x),s-x\rangle+\frac{\curv}{2}\gamma^{2}.
\tag{10}
\]
\end{lemma}
\begin{proof}
By definition of $\curv$ (Assumption~\ref{ass:curvature}) and rearranging the inequality.
\end{proof}

\begin{lemma}[Gap lower bound via pyramidal width]\label{lem:gap}
Let $g_k:=\gap_k=\langle\nabla f(x_k),x_k-s_k\rangle$.
Then
\[
g_k\;\ge\;\delta\;\sqrt{2L\,h_k}.
\tag{11}
\]
\end{lemma}
\begin{proof}
Consider the optimal solution $x^\star\in\C$. Convexity gives
\[
h_k = f(x_k)-f(x^\star)\le\langle\nabla f(x_k),x_k-x^\star\rangle .
\tag{12}
\]
Write $x^\star$ as a convex combination of vertices; among those vertices there exists $v\in S_k$ such that
\[
\langle\nabla f(x_k),x_k-v\rangle\ge
\langle\nabla f(x_k),x_k-s_k\rangle = g_k .
\tag{13}
\]
The definition of pyramidal width $\delta$ (Assumption~\ref{ass:polytope}) yields
\[
\frac{g_k}{\|\nabla f(x_k)\|\,\|x_k-s_k\|}
\;\ge\;\delta .
\tag{14}
\]
Since $\|x_k-s_k\|\le\diam(\C)$ and $\|\nabla f(x_k)\|\le\sqrt{2L\,h_k}$ (smoothness + convexity), we obtain \eqref{11}.
\end{proof}

\begin{lemma}[One‑step decrease]\label{lem:onestep}
Let $d_k$ be the direction chosen by AFW and $\gamma_k$ the exact line‑search (4). Then
\[
h_{k+1}\;\le\;h_k-\frac{g_k^{2}}{2L\,\diam(\C)^{2}} .
\tag{15}
\]
\end{lemma}
\begin{proof}
From Lemma~\ref{lem:descent} with $\gamma=\gamma_k$ and $s$ being either $s_k$ (FW) or $v_k$ (away):
\[
f(x_{k+1})\le f(x_k)+\gamma_k\langle\nabla f(x_k),d_k\rangle
+\frac{\curv}{2}\gamma_k^{2}.
\tag{16}
\]
Because $\gamma_k$ minimizes the RHS over $[0,\gamma_{\max}]$, the derivative at $\gamma_k$ (if interior) satisfies
\[
\langle\nabla f(x_k),d_k\rangle+\curv\gamma_k=0
\quad\Longrightarrow\quad
\gamma_k=-\frac{\langle\nabla f(x_k),d_k\rangle}{\curv}.
\tag{17}
\]
Plugging \eqref{17} into \eqref{16} gives
\[
f(x_{k+1})\le f(x_k)-\frac{\langle\nabla f(x_k),d_k\rangle^{2}}{2\curv}.
\tag{18}
\]
Since $\curv\le L\,\diam(\C)^{2}$ and $|\langle\nabla f(x_k),d_k\rangle|\ge g_k$ (by construction of $d_k$),
\[
h_{k+1}\le h_k-\frac{g_k^{2}}{2L\,\diam(\C)^{2}} .
\tag{19}
\]
\end{proof}

%%%%%%%%%%%%%%%%%%%%%%%%%%%%%%%%%%%%%%%%%%%%%%%%%%%%%%%%%%%%
\section*{Main Proof (Step‑by‑Step Derivation)}\label{sec:proof}
%%%%%%%%%%%%%%%%%%%%%%%%%%%%%%%%%%%%%%%%%%%%%%%%%%%%%%%%%%%%
We now combine Lemmas~\ref{lem:gap} and \ref{lem:onestep}.

\begin{enumerate}
\item[Step 1.] Apply Lemma~\ref{lem:onestep}:
\[
h_{k+1}\le h_k-\frac{g_k^{2}}{2L\,\diam(\C)^{2}} .
\tag{20}
\]
\item[Step 2.] Lower‑bound $g_k$ using Lemma~\ref{lem:gap}:
\[
g_k\ge\delta\sqrt{2L\,h_k}.
\tag{21}
\]
\item[Step 3.] Substitute \eqref{21} into \eqref{20}:
\[
h_{k+1}\le h_k-
\frac{\bigl(\delta\sqrt{2L\,h_k}\bigr)^{2}}{2L\,\diam(\C)^{2}}
= h_k-\frac{\delta^{2}}{L\,\diam(\C)^{2}}\,h_k .
\tag{22}
\]
\item[Step 4.] Define $\rho:=\dfrac{\delta^{2}}{L\,\diam(\C)^{2}}$ (cf.~\eqref{7}) and rewrite \eqref{22} as
\[
h_{k+1}\le (1-\rho)h_k .
\tag{23}
\]
\item[Step 5.] Recursively applying \eqref{23} yields
\[
h_k\le (1-\rho)^{k}h_0 .
\tag{24}
\]
\item[Step 6.] Using Lemma~\ref{lem:descent} with $\gamma=1$ gives
\[
g_k\le L\,\diam(\C)^{2}\,\gamma_k\le L\,\diam(\C)^{2},
\]
and together with \eqref{24} we obtain the bound on the duality gap
\[
g_k\le 2L\,\diam(\C)^{2}(1-\rho)^{k}h_0 .
\tag{25}
\]
\end{enumerate}
Equations \eqref{23}–\eqref{25} constitute the full derivation of Theorem~\ref{thm:linear}. $\square$

\section*{Rate Derivation (With Constants)}
From \eqref{23} we have a linear contraction factor $1-\rho$ with
\[
\rho=\frac{\delta^{2}}{L\,\diam(\C)^{2}}.
\]
If we denote $\kappa:=L\,\diam(\C)^{2}/\delta^{2}$, then $\rho=1/\kappa$ and
\[
h_k\le\Bigl(1-\frac{1}{\kappa}\Bigr)^{k}h_0
\le\exp\!\bigl(-k/\kappa\bigr)h_0 .
\tag{26}
\]
Thus AFW reaches an $\varepsilon$‑optimal solution after at most
\[
k\;\ge\;\kappa\log\!\bigl(h_0/\varepsilon\bigr)
\]
iterations.

\section*{Special Cases}
\begin{itemize}
\item \textbf{Simplex $\Delta_d$.} Here $\diam(\Delta_d)=\sqrt{2}$ and $\delta=2/\sqrt{d(d+1)}$. Hence
\[
\rho=\frac{4}{L\,d(d+1)} .
\]
The iteration complexity scales linearly with the dimension $d$.
\item \textbf{Line segment $[a,b]$.} $\diam=|b-a|$, $\delta=1$, giving $\rho=1/(L|b-a|^{2})$, which matches the exact contraction observed in the dry‑run.
\item \textbf{Strongly convex $f$.} If $f$ is $\mu$‑strongly convex, the duality gap satisfies $g_k\ge 2\mu h_k$, leading to an even larger contraction factor $\rho'=\min\{\rho,\,\mu/L\}$.
\end{itemize}

%%%%%%%%%%%%%%%%%%%%%%%%%%%%%%%%%%%%%%%%%%%%%%%%%%%%%%%%%%%%
\section*{References}
%%%%%%%%%%%%%%%%%%%%%%%%%%%%%%%%%%%%%%%%%%%%%%%%%%%%%%%%%%%%
\begin{thebibliography}{9}
\bibitem{lacoste2015linear}
\textsc{Lacoste-Julien, S.}, \textsc{Jaggi, M.} (2015). \emph{On the linear convergence of Frank–Wolfe variants}. 
Advances in Neural Information Processing Systems, 28, 1999–2007.
\end{thebibliography}

\end{document}